\subsection{Divisor classes in Krull domains}

Let A be a Krull domain. Recall that the group D(A) of divisors of A is the free commutative group generated by the set P(A) of its extremal elements (no. 3, Theorem 2) and that P(A) is identified with the set of prime ideals of A of height 1 (no. 6); for \(p \in P(A)\) we shall denote by \(v_{p}\) the normed essential valuation corresponding to p (no. 4); recall that the ring of \(v_{p}\) is \(A_{p}\) (no. 4, Corollary 1 to Proposition 6). We shall denote by F(A) the subgroup of D(A) consisting of the principal divisors and by C(A) = D(A)/F(A) the divisor class group of A (no. 2).

PROPOSITION 14. Let A be a Krull domain and B a Krull domain containing A. Suppose that the following condition holds:

(PDE) For every prime ideal \(\mathfrak{P}\) of \(\mathbf{B}\) of height 1, the prime ideal \(\mathfrak{P} \cap \mathbf{A}\) is zero or of height 1.

For \(\mathfrak{p} \in \mathrm{P}(\mathbf{A})\) the \(\mathfrak{P} \in \mathrm{P}(\mathbf{B})\) such that \(\mathfrak{P} \cap \mathbf{A} = \mathfrak{p}\) are finite in number; we write

\[
i (\mathfrak {p}) = \sum_ {\mathfrak {P} \in \mathrm{P} (\mathrm{B}), \mathfrak {P} \cap \mathrm{A} = \mathfrak {p}} e (\mathfrak {P} / \mathfrak {p}) \mathfrak {P},
\]

where \(e(\mathfrak{P} / \mathfrak{p})\) denotes the ramification index of \(v_{\mathfrak{P}}\) over \(v_{\mathfrak{p}}\) (Chapter VI, § 8, no. 1). Then \(i\) defines, by linearity, an increasing homomorphism of D(A) to D(B), which enjoys the following properties:

\begin{enumerate}
\def\labelenumi{(\alph{enumi})}
\tightlist
\item
  for every non-zero element \(x\) of the field of fractions of \(\mathbf{A}\) ,
\end{enumerate}

\[
i \left(\operatorname{div} _ {\mathbf {A}} (x)\right) = \operatorname{div} _ {\mathbf {B}} (x);
\]

\begin{enumerate}
\def\labelenumi{(\alph{enumi})}
\setcounter{enumi}{1}
\tightlist
\item
  for all D, D' in D(A),
\end{enumerate}

\[
i \left(\sup (\mathbf {D}, \mathbf {D} ^ {\prime})\right) = \sup (i (\mathbf {D}), i (\mathbf {D} ^ {\prime})).
\]

Let \(\mathfrak{p} \in \mathrm{P}(\mathbf{A})\) ; consider a non-zero element \(a\) of \(\mathfrak{p}\) ; the \(\mathfrak{P} \in \mathrm{P}(\mathbf{B})\) which contain \(a\) are finite in number (no. 6, Theorem 4); a fortiori the \(\mathfrak{P} \in \mathrm{P}(\mathbf{B})\) such that \(\mathfrak{P} \cap \mathbf{A} = \mathfrak{p}\) are finite in number.

We now show (a). By additivity, it may be assumed that \(x \in \mathbf{A}^* = \mathbf{A} - \{0\}\) . By definition, \(\operatorname{div}_{\mathbf{B}}(x) = \sum_{\mathfrak{P} \in \mathrm{P}(\mathbf{B})} v_{\mathfrak{P}}(x) \cdot \mathfrak{P}\). For all \(\mathfrak{P} \in \mathrm{P}(\mathbf{B})\) such that \(v_{\mathfrak{P}}(x) > 0\) , \(\mathfrak{P} \cap \mathbf{A}\) is non-zero (for \(x \in \mathfrak{P}\) ) and is therefore of height 1 by (PDE); setting \(\mathfrak{p} = \mathfrak{P} \cap \mathbf{A}\) , by definition of the ramification index, \(v_{\mathfrak{P}}(x) = e(\mathfrak{P}/\mathfrak{p})v_{\mathfrak{p}}(x)\) (since \(v_{\mathfrak{p}}\) and \(v_{\mathfrak{P}}\) are normed). As \(\operatorname{div}_{\mathbf{A}}(x) = \sum_{\mathfrak{p} \in \mathrm{P}(\mathbf{A})} v_{\mathfrak{p}}(x) \cdot \mathfrak{p}\), and \(i(\mathfrak{q}) = 0\) for all \(\mathfrak{q} \in \mathrm{P}(\mathbf{A})\) which is not of the form \(\mathfrak{Q} \cap \mathbf{A}\) where \(\mathfrak{Q} \in \mathrm{P}(\mathbf{B})\) , we deduce (a).

To prove (b) we write

\[
\mathrm{D} = \sum_ {\mathfrak {p} \in \mathrm{P} (\mathrm{A})} n (\mathfrak {p}) \cdot \mathfrak {p} \quad \text { and } \quad \mathrm{D} ^ {\prime} = \sum_ {\mathfrak {p} \in \mathrm{P} (\mathrm{A})} n ^ {\prime} (\mathfrak {p}) \cdot \mathfrak {p};
\]

the coefficient of \(\mathfrak{p}\) in \(\sup (\mathrm{D},\mathrm{D}^{\prime})\) is \(\sup (n(\mathfrak{p}),n^{\prime}(\mathfrak{p}))\) . Let \(\mathfrak{P}\) be an element of \(\mathrm{P(B)}\) . If \(\mathfrak{P}\cap \mathrm{A} = (0)\) , the coefficients of \(\mathfrak{P}\) in \(i(\mathrm{D})\) and \(i(\mathrm{D}^{\prime})\) and hence also in \(\sup (i(\mathrm{D}),i(\mathrm{D}^{\prime}))\) are zero; therefore the coefficient of \(\mathfrak{P}\) in \(i(\sup (\mathrm{D},\mathrm{D}^{\prime}))\) is zero. If \(\mathfrak{P}\cap \mathrm{A}\neq (0)\) , it is a prime ideal \(\mathfrak{p}\) of height 1 (by (PDE)); writing \(e = e(\mathfrak{P} / \mathfrak{p})\) , the coefficients of \(\mathfrak{P}\) in \(i(\mathrm{D})\) , \(i(\mathrm{D}^{\prime})\) and \(i(\sup (\mathrm{D},\mathrm{D}^{\prime}))\) are respectively \(en(\mathfrak{p}),en^{\prime}(\mathfrak{p})\) and \(e \cdot\sup (n(\mathfrak{p}),n^{\prime}(\mathfrak{p}))\) ; that of \(\sup (i(\mathrm{D}),i(\mathrm{D}^{\prime}))\) is

\[
\sup (e \cdot n (\mathfrak {p}), e \cdot n ^ {\prime} (\mathfrak {p})) = e \cdot \sup (n (\mathfrak {p}), n ^ {\prime} (\mathfrak {p})).
\]

This proves (b).

Under the hypotheses of Proposition 14, it follows from (a) that i defines, by taking quotients, a homomorphism \(\bar{i}\) , called canonical, of C(A) to C(B), which we shall also sometimes write as i, by an abuse of notation.

The condition (PDE) is fulfilled in the two following cases:

\begin{enumerate}
\def\labelenumi{(\arabic{enumi})}
\tightlist
\item
  B is integral over A; in this case, for the prime ideal P of B to be of height 1, it is necessary and sufficient that \(p = P \cap A\) be of height 1. For (0) is the only prime ideal of B lying above the ideal (0) of A (Chapter V, § 2, no. 1, Corollary 1 to Proposition 1); if P is of height 1, then \(p \neq 0\) ; if p were not of height 1, there would exist a prime ideal \(p'\) of A distinct from (0) and p and such that
\end{enumerate}

\(0 \subset p' \subset p\) ; but then, as B is an integral domain and A is an integrally closed domain, there would be a prime ideal \(P'\) of B such that \(P' \cap A = p'\) and \(P' \subset P\) (Chapter V, §2, no. 4, Theorem 3), contrary to the hypothesis. Conversely, if p is of height 1, there can exist no prime ideal \(P'\) of B distinct from 0 and P and such that \(0 \subset P' \subset P\) , otherwise \(0 \subset P' \cap A \subset p\) and \(P' \cap A\) would be distinct from 0 and p by virtue of Chapter V, §2, no. 1, Corollary 1 to Proposition 1.

\begin{enumerate}
\def\labelenumi{(\arabic{enumi})}
\setcounter{enumi}{1}
\tightlist
\item
  B is a flat A-module. More precisely:
\end{enumerate}

PROPOSITION 15. Let A and B be Krull domains such that B contains A and is a flat A-module. Then:

\begin{enumerate}
\def\labelenumi{(\alph{enumi})}
\item
  condition (PDE) of Proposition 14 is fulfilled;
\item
  for every divisorial ideal \(\mathfrak{a}\) of \(\mathbf{A}\) , \(\mathbf{B}\mathfrak{a}\) is the divisorial ideal of \(\mathbf{B}\) which corresponds to the divisor \(i(\mathrm{div}_{\mathbf{A}}(\mathfrak{a}))\) .
\end{enumerate}

To show (a), suppose that there exists a prime ideal \(\mathfrak{P}\) of B of height 1 such that \(\mathfrak{P} \cap A\) is neither 0 nor of height 1. Take an element \(x \neq 0\) in \(\mathfrak{P} \cap A\) . The ideals \(\mathfrak{p}_i\) of A of height 1 which contain \(x\) are finite in number and none contains \(\mathfrak{P} \cap A\) ; there therefore exists an element \(y\) of \(\mathfrak{P} \cap A\) such that \(y \notin \mathfrak{p}_i\) for all \(i\) (Chapter II, §2, no. 1, Proposition 2). Thus \(\operatorname{div}_{A}(x)\) and \(\operatorname{div}_{A}(y)\) are relatively prime elements of the ordered group P(A), so that \(\sup (\operatorname{div}_{A}(x), \operatorname{div}_{A}(y)) = \operatorname{div}_{A}(x) + \operatorname{div}_{A}(y) = \operatorname{div}_{A}(xy)\) ; as the ideals \(Ax \cap Ay\) and \(Axy\) are divisorial, we deduce that \(Ax \cap Ay = Axy\) . Since B is a flat A-module, \(Bx \cap By = Bxy\) (Chapter I, §2, no. 6, Proposition 6). This implies that \(\sup (v_{\mathfrak{P}}(x), v_{\mathfrak{P}}(y)) = v_{\mathfrak{P}}(xy) = v_{\mathfrak{P}}(x) + v_{\mathfrak{P}}(y)\) , which contradicts the inequalities \(v_{\mathfrak{P}}(x) > 0\) , \(v_{\mathfrak{P}}(y) > 0\) (which hold since \(x\) and \(y\) are in \(\mathfrak{P}\) ). Thus (a) has been proved by reductio ad absurdum.

We now show (b). If \(a\) is a divisorial ideal of \(A\) , it is the intersection of two fractional principal ideals (no. 5, Corollary 2 to Proposition 9), say

\[
a = d ^ {- 1} (\mathrm{A} a \cap \mathrm{A} b)
\]

where \(a, b, d\) are in \(\mathbf{A}^*\) ; as \(\mathbf{B}\) is flat over \(\mathbf{A}, \mathbf{B}a = d^{-1}(\mathbf{B}a \cap \mathbf{B}b)\) (Chapter I, § 2, no. 6, Proposition 6), which shows that \(\mathbf{B}a\) is divisorial. This shows also that \(\mathrm{div}_{\mathbf{B}}(\mathbf{B}a) = \sup (\mathrm{div}_{\mathbf{B}}(a), \mathrm{div}_{\mathbf{B}}(b)) - \mathrm{div}_{\mathbf{B}}(d)\) ; using Proposition 14 (a) and (b), it is seen that

\[
\begin{array}{r l} \operatorname{div} _ {\mathbf {B}} (\mathrm{Ba}) & = \sup (i (\operatorname{div} _ {\mathbf {A}} (a)), i (\operatorname{div} _ {\mathbf {A}} (b))) - i (\operatorname{div} _ {\mathbf {A}} (d)) \\ & = i (\sup (\operatorname{div} _ {\mathbf {A}} (a), \operatorname{div} _ {\mathbf {A}} (b))) - i (\operatorname{div} _ {\mathbf {A}} (d)) \\ & = i (\operatorname{div} _ {\mathbf {A}} (\mathrm{A} a \cap \mathrm{Ab})) - i (\operatorname{div} _ {\mathbf {A}} (d)) \\ & = i (\operatorname{div} _ {\mathbf {A}} (d ^ {- 1} (\mathrm{A} a \cap \mathrm{Ab}))) = i (\operatorname{div} _ {\mathbf {A}} (\mathrm{a})). \end{array}
\]

COROLLARY. Let A be a local Krull domain and B a discrete valuation ring such that B dominates A and is a flat A-module. Then A is a field or a discrete valuation ring.

Let M be the maximal ideal of B. By (PDE), \(M \cap A\) is zero or of height 1. As it is, by hypothesis the maximal ideal of A, our assertion follows from Proposition 11 of no. 7.

Remark. In the former of the two above cases, the mapping \(i\colon\mathrm{D}(\mathrm{A})\to\mathrm{D}(\mathrm{B})\) is injective: as the elements of \(\mathrm{P}(\mathrm{B})\) form a basis of \(\mathrm{D}(\mathrm{B})\) and two distinct ideals of \(\mathrm{P}(\mathrm{A})\) cannot be the traces on \(\mathrm{A}\) of the same ideal of \(\mathrm{P}(\mathrm{B})\) , it amounts to verifying that \(i(\mathfrak{p})\neq0\) for all \(\mathfrak{p}\in\mathrm{P}(\mathrm{A})\) ; now, this follows from Chapter V, § 2, no. 1, Theorem 1. It is similarly seen that, if \(\mathrm{B}\) is a faithfully flat A-module, \(i\) is injective (Chapter II, § 2, no. 5, Corollary 4 to Proposition 11).

In what follows, we propose to study the canonical homomorphism \(\bar{\imath}\) from C(A) to C(B) for certain ordered pairs of Krull domains A, B.

PROPOSITION 16. Let A be a Zariski ring such that its completion \(\hat{\mathbf{A}}\) is a Krull domain. Then A is a Krull domain and the canonical homomorphism \(\bar{i}\) from C(A) to C( \(\hat{\mathbf{A}}\) ) (which is defined since \(\hat{\mathbf{A}}\) is a flat A-module; cf.~Chapter III, § 3, no. 4, Theorem 3) is injective.

As \(\hat{\mathbf{A}}\) is an integral domain and \(\mathbf{A} \subset \hat{\mathbf{A}}\) , \(\mathbf{A}\) is an integral domain. Let \(\mathbf{L}\) be the field of fractions of \(\hat{\mathbf{A}}\) and \(\mathbf{K} \subset \mathbf{L}\) that of \(\mathbf{A}\) . As \(\mathbf{A} = \hat{\mathbf{A}} \cap \mathbf{K}\) (Chapter III, § 3, no. 5, Corollary 4 to Proposition 9), \(\mathbf{A}\) is a Krull domain (no. 3, Example 4). The fact that \(\bar{i}: \mathbf{C}(\mathbf{A}) \to \mathbf{C}(\hat{\mathbf{A}})\) is injective follows from Proposition 15 (b) and the fact that, if \(\mathfrak{b}\hat{\mathbf{A}}\) is principal, \(\mathfrak{b}\) is principal (Chapter III, § 3, no. 5, Corollary 3 to Proposition 9).

Now let A be a Krull domain and S a multiplicative subset of A not containing 0. The group D(A) (resp. D(S \(^{-1}\) A) is the free commutative group with basis the set of div(p) (resp. div(S \(^{-1}\) p)), where p runs through the set of prime ideals of A of height 1 (resp. the set of prime ideals of A of height 1 such that p ∩ S = ∅) (no. 4, Proposition 6) and, if p ∩ S = ∅, then i(div(p)) = div(S \(^{-1}\) p). Thus D(S \(^{-1}\) A) is identified with the direct factor of D(A) generated by the elements div(p) such that p ∩ S = ∅ and admits as complement the free commutative subgroup of D(A) with basis the set of div(p) such that p ∩ S ≠ ∅; we shall denote this complement by G. As i: D(A) → D(S \(^{-1}\) A) is surjective, so is i: C(A) → C(S \(^{-1}\) A); and:

\[
\mathbf {G} / (\mathbf {G} \cap \mathbf {F} (\mathbf {A})) = (\mathbf {G} + \mathbf {F} (\mathbf {A})) / \mathbf {F} (\mathbf {A}) = \operatorname{Ker} (\bar {i});\tag{5}
\]

for if an element of \(\mathrm{D}(\mathrm{S}^{-1}\mathrm{A})\) is equal to \(\operatorname{div}_{\mathbf{S}^{-1}\mathbf{A}}(x / s)\) , where \(x \in \mathbf{A}\) and \(s \in \mathbf{S}\) , it is the image under \(i\) of the principal divisor \(\operatorname{div}_{\mathbf{A}}(x)\) (Proposition 14).

Suppose now that S is generated by a family of elements \((p_{i})_{i\in I}\) of A such that the principal ideals \(Ap_{i}\) are all prime. Then, if \(\mathfrak{p}\) is a prime ideal of A of height 1 such that \(\mathfrak{p} \cap \mathrm{S} \neq \varnothing\) , \(\mathfrak{p}\) contains a product of powers of the \(p_{i}\) and therefore one of the \(p_{i}\) , say \(p_{\alpha}\) ; as \(\mathrm{Ap}_{\alpha}\) is non-zero and prime and \(\mathfrak{p}\) is of height 1, it follows that \(\mathfrak{p} = \mathrm{Ap}_{\alpha}\) . In the above notation, therefore \(\mathrm{G} \subset \mathrm{F}(\mathrm{A})\) and (5) shows that the kernel of \(\bar{i}\) is zero. We have therefore shown the following result:

PROPOSITION 17. Let A be a Krull domain and S a multiplicative subset of A not containing 0. Then the canonical homomorphism \(\bar{i}\) from C(A) to C(S \(^{-1}\) A) is surjective. If further S is generated by a family of elements \(p_{i}\) such that the principal ideals Ap \(_{i}\) are all prime, then \(\bar{i}\) is bijective.

As a second application of formula (5), consider the following situation: let R be a Krull domain; take A to be the polynomial ring A = R{[}X{]} (no. 9, Proposition 13) and S to be the set R - (0) of non-zero constant polynomials of A. The prime ideals p of A of height 1 such that \(p \cap S \neq \varnothing\) are those of the form \(p_{0}A\) , where \(p_{0}\) is a prime ideal of R of height 1 (no. 9, Remark). Hence, in the notation introduced above, G is identified with D(R) by identifying \(\operatorname{div}_{\mathbf{A}}(\mathfrak{p}_{0}\mathbf{A})\) with \(\operatorname{div}_{\mathbf{R}}(\mathfrak{p}_{0})\) . On the other hand \(\mathbf{G} \cap \mathbf{F}(\mathbf{A})\) is identified with \(\mathbf{F}(\mathbf{R})\) : for if an ideal \(a_{0}\) of R generates a principal ideal \(a_{0}\mathbf{A} = f(\mathbf{X})\mathbf{A}\) in A = R{[}X{]}, then \(f(0) \in a_{0}\mathbf{A}\) since \(a_{0}\mathbf{A}\) is a graded ideal of the ring A (graded by the usual degree of polynomials) and hence \(f(0) \in a_{0}\) ; further, for \(a \in a_{0}\) , \(a = f(\mathbf{X})g(\mathbf{X})\) where \(g(\mathbf{X}) \in \mathbf{R}\) , whence, comparing terms of degree 0, \(a = f(0)g(0)\) ; it follows that \(a_{0}\) is the principal ideal of R generated by \(f(0)\) . Finally, denoting by K the field of fractions of R, \(S^{-1}A\) is identified with the polynomials ring K{[}X{]}, which is a principal ideal domain; hence \(\mathrm{C}(\mathrm{S}^{-1}\mathrm{A}) = (0)\) . Thus, by (5), \(\mathrm{C}(\mathrm{A}) = \mathrm{Ker}(\bar{i})\) is identified with C(R) and we have proved the following result:

PROPOSITION 18. Let R be a Krull domain and A the polynomial ring R{[}X{]}. The canonical homomorphism of C(R) to C(R{[}X{]}) is bijective.

